% ----- CHAUPAI 16 -----
\subsection*{\deva{षोडश चौपाई} \(Sixteenth Chaupai\)}
\addcontentsline{toc}{subsection}{\deva{षोडश चौपाई} \(Sixteenth Chaupai\) — \deva{तुम उपकार...}}


\begin{minipage}[t]{0.55\textwidth}
\vspace{0pt}
{\large \linespread{1.5}\selectfont
\deva{तुम उपकार सुग्रीवहिं कीन्हा।}\\
\deva{राम मिलाय राज पद दीन्हा॥}
\par}

\vspace{0.5em}

{\large \linespread{1.5}\selectfont
\telu{తుమ ఉపకార సుగ్రీవహిఁ కీన్హా।}\\
\telu{రామ మిలాయ రాజ పద దీన్హా॥}
\par}

\vspace{0.5em}

{\large \linespread{1.3}\selectfont
\textit{tuma upakāra sugrīvahĩ kīnhā |}\\
\textit{rāma milāya rāja pada dīnhā ||}
\par}
\end{minipage}%
\hfill
\begin{minipage}[t]{0.42\textwidth}
\vspace{0pt}
\begin{tcolorbox}[anchorbox]
\textbf{Lifting Others Up:} Sugriva was a displaced outcast who had lost everything. Hanuman didn't just use his skills to advance his own career; he connected Sugriva to Rama and helped restore him to a position of leadership. A truly great leader uses their network and talents to empower and uplift those who are struggling.
\vspace{0.5em}\newline Hanuman leveraged his own networking and diplomatic skills to uplift a struggling, displaced friend into a position of leadership.\newline\null\hfill\href{https://www.valmiki.iitk.ac.in/sloka?field_kanda_tid=4&language=dv&field_sarga_value=5}{\textit{Kishkindha Kanda, 5:15-18}}
\end{tcolorbox}
\end{minipage}


\vspace{0.5em}
\subsubsection*{\textbf{Visual Rhythm Map:}}
\begin{table}[h!]
\centering
\renewcommand{\arraystretch}{1.2}
\setlength{\tabcolsep}{0pt}
\begin{tabularx}{\textwidth}{|*{16}{>{\centering\arraybackslash\hsize=1.0\hsize}X|}}
\hline
\multicolumn{2}{|>{\centering\arraybackslash\hsize=2.0\hsize}X|}{\textbf{\laghu{तु}\laghu{म}} \par (2)} &
\multicolumn{5}{>{\centering\arraybackslash\hsize=5.0\hsize}X|}{\textbf{\laghu{उ}\laghu{प}\guru{का}\laghu{र}} \par (5)} &
\multicolumn{5}{>{\centering\arraybackslash\hsize=5.0\hsize}X|}{\textbf{\laghu{सु}\guru{ग्री}\laghu{व}\laghu{हिं}} \par (5)} &
\multicolumn{4}{>{\centering\arraybackslash\hsize=4.0\hsize}X|}{\textbf{\guru{की}\guru{न्हा}} \par (4)} \\
\hline
\end{tabularx}
\vspace{-1.5em}
\end{table}

\begin{table}[h!]
\centering
\renewcommand{\arraystretch}{1.2}
\setlength{\tabcolsep}{0pt}
\begin{tabularx}{\textwidth}{|*{16}{>{\centering\arraybackslash\hsize=1.0\hsize}X|}}
\hline
\multicolumn{3}{|>{\centering\arraybackslash\hsize=3.0\hsize}X|}{\textbf{\guru{रा}\laghu{म}} \par (3)} &
\multicolumn{4}{>{\centering\arraybackslash\hsize=4.0\hsize}X|}{\textbf{\laghu{मि}\guru{ला}\laghu{य}} \par (4)} &
\multicolumn{3}{>{\centering\arraybackslash\hsize=3.0\hsize}X|}{\textbf{\guru{रा}\laghu{ज}} \par (3)} &
\multicolumn{2}{>{\centering\arraybackslash\hsize=2.0\hsize}X|}{\textbf{\laghu{प}\laghu{द}} \par (2)} &
\multicolumn{4}{>{\centering\arraybackslash\hsize=4.0\hsize}X|}{\textbf{\guru{दी}\guru{न्हा}} \par (4)} \\
\hline
\end{tabularx}
\vspace{-1.5em}
\end{table}

\vspace{0.5em}
\textbf{ \deva{सरल-संस्कृतम्}:}\\
 \deva{भवान् सुग्रीवाय महान्तम् उपकारं कृतवान्।}\\
 \deva{श्रीरामचन्द्रेण सह तस्य मेलनं कारयित्वा, तस्मै राजपदं दत्तवान्।}\\

You did a great favor for Sugriva.\\
By uniting him with Lord Rama, you bestowed the royal throne upon him.


\vspace{1em}
\newpage
\textbf{\deva{प्रतिपदार्थः} (Meaning of each word):}
\begin{table}[h!]
\centering
\renewcommand{\arraystretch}{1.2}
\begin{tabularx}{\textwidth}{@{} l l X X @{}}
\toprule
\textbf{अवधी} & \textbf{संस्कृतम्} & \textbf{English} & \textbf{Grammar \& Notes} \\
\midrule
\deva{तुम} \gotchaicon / \telu{తుమ} / \textit{tuma}  &  \deva{त्वम्}  &  You  &   \\
\deva{उपकार} / \telu{ఉపకార} / \textit{upakāra}  &  \deva{उपकारम्}  &  Favor / Service  &   \\
\deva{सुग्रीवहिं} \gotchaicon / \telu{సుగ్రీవహిఁ} / \textit{sugrīvahĩ}  &  \deva{सुग्रीवाय}  &  To Sugriva  & Objective suffix `-hiṃ` \\
\deva{कीन्हा} \gotchaicon / \telu{కీన్హా} / \textit{kīnhā}  &  \deva{कृतवान्}  &  Did  & Awadhi past tense \\
\deva{राम मिलाय} / \telu{రామ మిలాయ} / \textit{rāma milāya}  &  \deva{रामेण मेलयित्वा}  &  Having united with Rama  &   \\
\deva{राज पद} / \telu{రాజ పద} / \textit{rāja pada}  &  \deva{राज-पदम्}  &  Royal Throne  &   \\
\deva{दीन्हा} / \telu{దీన్హా} / \textit{dīnhā}  &  \deva{दत्तवान्}  &  Gave  & Rhymes with keenhaa \\
\bottomrule
\end{tabularx}
\end{table}

\begin{tcolorbox}[poetrybox]
\textbf{\deva{त्वम्} $\rightarrow$ \deva{तुम} \gotchaicon\  Pronoun Across Millennia:}\\
Sanskrit uses \deva{त्वम्} for ''you'' (nominative singular) that Awadhi collapses into the simple, \deva{तुम}: the exact same form that survived into modern Hindi, Urdu, and Bengali. 
Note also that \deva{तुम} in Awadhi works for all cases;
 there is no separate accusative or dative form.\\[0.4em]
\textbf{\deva{कृतवान्} $\rightarrow$ \deva{कीन्हा} \gotchaicon\  The Awadhi Past Tense Signature:}\\
\deva{कीन्हा} is the canonical Awadhi perfect past form of \deva{करना} (to do). 
Notice the dramatic reduction: Sanskrit requires an elaborate perfect participle \deva{कृतवान्} (''one who has done''), while Awadhi achieves the same meaning with a single crisp word. 
The rhyming pair \deva{कीन्हा} / \deva{दीन्हा} (did / gave) is a beautiful structural device Tulsidas uses to link this couplet's two actions musically.\\[0.4em]
\textbf{\deva{सुग्रीवहिं} — A Metrical Flex:}\\
\deva{सुग्रीवहिं} carries 5 matras in the rhythm map (\deva{सु}{\scriptsize(1)} + \deva{ग्री}{\scriptsize(2)} + \deva{व}{\scriptsize(1)} + \deva{हिं}{\scriptsize(2)} = 5), yet a phonologically valid 6-matra reading also exists (\deva{सु}{\scriptsize(2)} + \deva{ग्री}{\scriptsize(2)} + \deva{व}{\scriptsize(1)} + \deva{हि}{\scriptsize(1)} = 6). Compare this with \deva{पुत्र} from Chaupai 2 — both words contain a conjunct consonant (\deva{संयुक्ताक्षर}), yet they are weighted differently. The distinction rests on a classical \textit{Apavāda Sūtra} (exception rule) governing stress (\textit{Pratibhā}):

\begin{itemize}
  \item \textbf{\deva{पुत्र} (pu-tra):} The conjunct \deva{त्र} (t+ra) involves a \textit{Sparśa} (stop) consonant. 
  The hard stop forces weight onto the preceding \deva{पु}, which becomes \textit{Guru} by the standard \textit{Saṃyuktapūrva} rule. 
  Result: \textit{Guru}+\textit{Laghu} = 3 matras.

  \item \textbf{\deva{सुग्री} (su-grī):} The conjunct \deva{ग्री} involves a soft consonant followed by \deva{र} (Repha). 
  Sanskrit and Telugu prosody permit a \textit{Śithila} (soft/lax) exception here: when the cluster is pronounced lightly, the preceding short vowel need \emph{not} become \textit{Guru}, keeping \deva{सु} as \textit{Laghu}. 
  Result: \textit{Laghu}+\textit{Guru} = 3 matras, but the opening syllable is light.
\end{itemize}

In the 16-matra Chaupai framework, Tulsidas invokes this \textit{Śithila} license to hold \deva{सुग्रीवहिं} at 5 matras, fitting the meter perfectly. 
The 6-matra reading remains phonologically available for singers who choose to lean into the conjunct. 
This is \textit{Bhāva Pradhāna} prosody at its finest: the emotional weight of the meaning guides the syllabic weight of the delivery.
\end{tcolorbox}

\newpage
